\documentclass[14pt]{extarticle}
\usepackage{fontspec}
\usepackage{polyglossia}
\usepackage{amsmath}
\usepackage{amssymb}
\usepackage{xcolor}
\usepackage{fancyhdr}
\usepackage{graphicx}
\usepackage{listings}
\usepackage[most]{tcolorbox}
% \usepackage{setspace}
\usepackage{pifont}
\usepackage{enumitem}
\usepackage{titlesec}
\usepackage[bottom]{footmisc}
\usepackage{titling}
\usepackage{minted}
\usepackage{etoolbox}
\usepackage{array}
\usepackage{longtable} % page-breakable tables (must load before polyglossia/bidi RTL setup)
\usepackage{extsizes}
\usepackage{float}
\usepackage{booktabs}
\usepackage{hyperref}
\usepackage[onehalfspacing]{setspace}
\newif\ifwhiteblack
\ifdefined\setwhiteblack
  \setwhiteblack
\fi

\ifwhiteblack
  \usepackage[a4paper, landscape, top=1.25cm, bottom=1.25cm, left=1.25cm, right=1.25cm, headheight=1.5cm, headsep=0.5cm]{geometry}
\else
  \usepackage[a4paper, top=2.7cm, bottom=2.7cm, left=2.5cm, right=2.5cm, headheight=1.5cm, headsep=0.5cm]{geometry}
\fi
\usepackage{tikz}
\usepackage{datetime2}

\usetikzlibrary{automata, positioning, arrows.meta, arrows}
\providecommand{\LTcaptype}{table}

\newfontfamily\emoji{Segoe UI Emoji}

\pagestyle{fancy}

\setmainlanguage[numerals=western]{arabic}
\setotherlanguage{english}
\newfontfamily\arabicfont[Script=Arabic]{Calibri}
\newfontfamily\arabicfonttt[Script=Arabic]{Courier New}



\ifwhiteblack
\lstset{
  language=[Sharp]C,
  numbers=left,
  stepnumber=1,
  numbersep=8pt,
  frame=single,
  basicstyle=\ttfamily\small,
  keywordstyle=\bfseries,
  stringstyle=\ttfamily,
  commentstyle=\itshape
}
\else
\lstset{
  language=[Sharp]C,
  numbers=left,
  stepnumber=1,
  numbersep=8pt,
  frame=single,
  basicstyle=\ttfamily\small,
  keywordstyle=\color{blue},
  stringstyle=\color{red},
  commentstyle=\color{green!50!black}
}
\fi

\newif\ifdetailed
\ifdefined\setdetailed
  \setdetailed
\fi

\newif\ifwithsols
\ifdefined\setwithsols
  \setwithsols
\fi

\ifwhiteblack
  \AtBeginDocument{
    \LARGE
    \usemintedstyle{bw}
  }
\fi

% Default minted look: compact code that wraps long lines.
% Not applied in black-and-white landscape mode, which already uses \LARGE.
\ifwhiteblack\else
  \setminted{fontsize=\small, breaklines=true}
\fi

% unified tcolorboxes for articles
\ifwhiteblack
\tcbset{colback=white, colframe=black, fonttitle=\bfseries, boxrule=1pt}
\newtcolorbox{boxDef}[1][]{title={{\emoji📘} تعريف\ifx\\#1\\\else ~#1\fi :}}
\newtcolorbox{boxExercise}[1][]{title={{\emoji🧩} تمرين\ifx\\#1\\\else ~#1\fi :}}
\newtcolorbox{boxExample}[1][]{title={{\emoji📝} مثال\ifx\\#1\\\else ~#1\fi :}}
\newtcolorbox{boxNote}[1][]{title={{\emoji✨} ملاحظة\ifx\\#1\\\else ~#1\fi :}}
\newtcolorbox{boxAttention}[1][]{title={{\emoji🔔} تنبيه\ifx\\#1\\\else ~#1\fi :}}
\newtcolorbox{boxWarning}[1][]{title={{\emoji⚡} ملاحظة هامة\ifx\\#1\\\else ~#1\fi :}}
\newtcolorbox{boxSolution}[1][]{title={{\emoji✅} حل\ifx\\#1\\\else ~#1\fi :}}
\newtcolorbox{boxSymbol}[1][]{title={{\emoji🔣} رمز\ifx\\#1\\\else ~#1\fi :}}
\newtcolorbox{boxHint}[1][]{title={{\emoji💡} تلميح\ifx\\#1\\\else ~#1\fi :}}
\else
\tcbset{colback=white, colframe=black, fonttitle=\bfseries, boxrule=0.8pt}
\newtcolorbox{boxDef}[1][]{colback=blue!5!white,colframe=blue!75!black,
  title={{\emoji📘} تعريف\ifx\\#1\\\else ~#1\fi :}}
\newtcolorbox{boxExercise}[1][]{colback=cyan!5!white,colframe=cyan!70!black,
  title={{\emoji🧩} تمرين\ifx\\#1\\\else ~#1\fi :}}
\newtcolorbox{boxExample}[1][]{colback=yellow!5!white,colframe=orange!90!black,
  title={{\emoji📝} مثال\ifx\\#1\\\else ~#1\fi :}}
\newtcolorbox{boxNote}[1][]{colback=gray!10!white,colframe=black,
  title={{\emoji✨} ملاحظة\ifx\\#1\\\else ~#1\fi :}}
\newtcolorbox{boxAttention}[1][]{colback=magenta!10!white,colframe=magenta!80!black,
  title={{\emoji🔔} تنبيه\ifx\\#1\\\else ~#1\fi :}}
\newtcolorbox{boxWarning}[1][]{colback=red!5!white,colframe=red!75!black,
  title={{\emoji⚡} ملاحظة هامة\ifx\\#1\\\else ~#1\fi :}}
\newtcolorbox{boxSolution}[1][]{colback=green!5!white,colframe=green!60!black,
  title={{\emoji✅} حل\ifx\\#1\\\else ~#1\fi :}}
\newtcolorbox{boxSymbol}[1][]{colback=purple!5!white,colframe=purple!70!black,
  title={{\emoji🔣} رمز\ifx\\#1\\\else ~#1\fi :}}
\newtcolorbox{boxHint}[1][]{colback=teal!5!white,colframe=teal!60!black,
  title={{\emoji💡} تلميح\ifx\\#1\\\else ~#1\fi :}}
\fi


\ifwhiteblack
\tcbset{simplecode/.style={ colback=white, colframe=black, boxrule=1pt, arc=2pt, left=4pt,right=4pt,top=4pt,bottom=4pt}}
\else
\tcbset{simplecode/.style={ colback=gray!5, colframe=black!50, boxrule=0.4pt, arc=2pt, left=4pt,right=4pt,top=4pt,bottom=4pt}}
\fi
% Compact code boxes: tighter line spacing (the preamble uses onehalfspacing) and small vertical gaps.
% Must come AFTER the simplecode definition above.
\tcbset{simplecode/.append style={fontupper=\setstretch{0.95}, top=4pt, bottom=4pt, before skip=6pt, after skip=6pt}}
\newenvironment{boxCode}{\begin{tcolorbox}[simplecode]}{\end{tcolorbox}}

% Code cell used inside tables (keeps code left-to-right inside RTL tables).
% Use inside a fixed-width p{...} column: \begin{pycell}\begin{minted}...\end{minted}\end{pycell}
\newtcolorbox{pycell}{frame empty,colback=white,boxsep=0pt,left=2pt,right=2pt,top=2pt,bottom=2pt,width=\linewidth}

\newcolumntype{C}[1]{>{\centering\arraybackslash}p{#1}}

\makeatletter
\DTMnewdatestyle{dotted}{%
  \renewcommand{\DTMdisplaydate}[4]{\two@digits{##3}.\two@digits{##2}.##1}%
  \renewcommand{\DTMDisplaydate}{\DTMdisplaydate}%
}
\makeatother
\DTMsetdatestyle{dotted}

% redefine spaces after titles
\makeatletter
\renewcommand{\@maketitle}{%
  \begin{center}
    {\huge \bfseries \@title \par}%
    \vskip 0.2em % space between title and author
    {\large \@author \par}%
    % \vskip 0.2em % space between author and date
    % {\normalsize \@date \par}%
    {\small \hfill \begin{english}\DTMtoday\end{english} \par}%
  \end{center}
}
\makeatother

\fancyhf{} % clear default
\fancypagestyle{plain}{
  \fancyhf{}
  \fancyhead[L]{مدرسة التسامح الشاملة}
  % \fancyhead[L]{\includegraphics[height=1cm]{../../../images/logoTasamoh.png}}
  \fancyhead[R]{الأستاذ محمود اغبارية}
  \fancyfoot[C]{\thepage}
}

\fancyhead[L]{مدرسة التسامح الشاملة}
\fancyhead[R]{الأستاذ محمود اغبارية}
\fancyfoot[C]{\thepage}

\setcounter{tocdepth}{3} % only section subsection and subsubsection in TOC

\newcommand{\importpdfpage}[2]{%
    \thispagestyle{empty}
    \newgeometry{left=0cm,right=0cm,top=0cm,bottom=0cm}
    \includegraphics[page=#2,width=0.95\paperwidth,height=0.95\paperheight,keepaspectratio]{#1}
    \restoregeometry
}

\newcommand{\insertFullImg}[1]{%
    \noindent
    \makebox[\textwidth][c]{
        \includegraphics[width=0.9\paperwidth,keepaspectratio]{#1}%
    }%
}%


\setminted{fontsize=\footnotesize, breaklines=true}
\tcbset{simplecode/.append style={fontupper=\setstretch{0.95}, top=4pt, bottom=4pt, before skip=6pt, after skip=6pt}}

\title{ملخص شامل للطلاب: البرمجة كائنية التوجه في \textenglish{C\#}}
\author{أ. محمود اغبارية}

\begin{document}

\maketitle
\thispagestyle{fancy}
\vspace{-0.6cm}

\begin{boxNote}[- دليل استخدام الملخص]
لكل موضوع في هذا الملخص: \textbf{شرح مختصر}، \textbf{مثال موضعي} قصير ومستقل مع كيفية استدعائه من \textenglish{Main}، و\textbf{تنبيهات} لأخطاء شائعة يجب الحذر منها أثناء الحل.
\end{boxNote}

% ==============================================================
\section*{الجزء الأول: أساسيات الفئات والتغليف}
% ==============================================================

\subsection*{1. الفئة والصفات والتغليف (عمليات \textenglish{Get} و \textenglish{Set})}
\textbf{الفكرة:} الفئة هي قالب لإنشاء كائنات. نُعرّف الصفات بإذن وصول \textenglish{private} لحمايتها من التعديل العشوائي من خارج الفئة، ونتيح قراءتها عبر عمليات \textenglish{Get} وتعديلها بأمان عبر عمليات \textenglish{Set} بعد فحص صحة القيمة المُتلقّاة.

\begin{boxCode}
\begin{english}
\begin{minted}{csharp}
public class Player
{
    private string name;
    private int score;
    public string GetName() { return this.name; }
    public void SetName(string name) { this.name = name; }
    public int GetScore() { return this.score; }
    public void SetScore(int score)
    {
        if (score >= 0) { this.score = score; } // Reject negative values
    }
}
// Usage in Main:
Player p = new Player();
p.SetName("Lina");
p.SetScore(50);
p.SetScore(-7); // Ignored because -7 is negative
Console.WriteLine(p.GetName() + ": " + p.GetScore()); // Prints: Lina: 50
\end{minted}
\end{english}
\end{boxCode}

\textbf{تنبيه:} لا تجعلوا الصفات \textenglish{public} للتسهيل؛ لأن ذلك يُفقد الفئة القدرة على التحقق من القيم. استعملوا \textenglish{this.score = score;} للتمييز بين صفة الكائن والمُعطى.

\filbreak
\subsection*{2. العمليات البنائية والتحميل الزائد (\textenglish{this})}
\textbf{الفكرة:} العملية البنائية تهيّئ صفات الكائن لحظة إنشائه بواسطة \textenglish{new}؛ وهي تحمل نفس اسم الفئة وليس لها نمط إرجاع (حتى \textenglish{void} لا تُكتب). يمكن كتابة عدة عمليات بنائية بمعطيات مختلفة (تحميل زائد)، ولتجنب تكرار الكود تستدعي إحداها الأخرى عبر \textenglish{: this(...)}.

\begin{boxCode}
\begin{english}
\begin{minted}{csharp}
public class Rect
{
    private int width;
    private int height;

    public Rect(int width, int height)
    {
        this.width = width;
        this.height = height;
    }

    // Square constructor calls the first constructor
    public Rect(int side) : this(side, side) { }
}

// Usage in Main:
Rect r1 = new Rect(4, 7); // Calls 1st constructor (4x7)
Rect r2 = new Rect(5);    // Calls 2nd constructor (5x5)
\end{minted}
\end{english}
\end{boxCode}

\textbf{تنبيه:} إذا كتبتم أي عملية بنائية في الفئة، تلغى العملية البنائية الافتراضية الفارغة \textenglish{new Rect()} تلقائياً؛ فإن احتجتموها يجب كتابتها صراحةً.

\filbreak
\subsection*{3. العملية \textenglish{ToString()} وتنسيق السلاسل النصية}
\textbf{الفكرة:} كل كائن يمتلك العملية \textenglish{ToString()}. نعيد تعريفها (\textenglish{override}) لتعيد نصاً يصف قيم صفات الكائن. عند تمرير الكائن إلى \textenglish{Console.WriteLine(obj)} تُستدعى \textenglish{ToString()} تلقائياً.

\begin{boxCode}
\begin{english}
\begin{minted}{csharp}
public class Player
{
    private string name;
    private int score;

    public override string ToString()
    {
        return $"Player:{this.name}, Score:{this.score}";
    }
}

// Usage in Main:
Console.WriteLine(p);       // Automatically calls p.ToString()
string info = p.ToString(); // Stores the returned string
\end{minted}
\end{english}
\end{boxCode}

\textbf{تنبيه:} نسيان الكلمة \textenglish{override} يجعل طباعة الكائن تطبع اسم الفئة فقط. انتبهوا أنّ \textenglish{ToString()} \textbf{تعيد} نصاً (\textenglish{return}) ولا تطبع داخلها بواسطة \textenglish{Console.WriteLine}.

\filbreak
\subsection*{4. عمليات تعيد \textenglish{bool} وفحص المعطيات}
\textbf{الفكرة:} العملية التي قد تفشل (مثل سحب مبلغ أكبر من الرصيد أو تمرير قيمة سالبة) تعيد \textenglish{true} عند النجاح و \textenglish{false} عند الفشل. نفحص الشروط \textbf{أولاً}، ولا نعدّل صفات الكائن إلا إذا تحققت جميع الشروط.

\begin{boxCode}
\begin{english}
\begin{minted}{csharp}
public class Wallet
{
    private int coins;

    public bool Spend(int amount)
    {
        if (amount <= 0 || amount > this.coins)
        {
            return false;
        }
        this.coins -= amount; // Modify state only after all checks pass
        return true;
    }
}

// Usage in Main:
if (w.Spend(30)) { Console.WriteLine("Success"); }
else { Console.WriteLine("Failed"); }
\end{minted}
\end{english}
\end{boxCode}

\filbreak
\subsection*{5. كائن من فئة أخرى كصفة (المؤشرات والعملية البنائية الناسخة)}
\textbf{الفكرة:} المتغير من نمط فئة يحمل \textbf{مؤشراً (عنواناً)} للكائن في الذاكرة وليس الكائن نفسه. كتابة \textenglish{b = a} تنسخ المؤشر فيشير المتغيران لنفس الكائن. لحماية صفات الكائن المركب من التعديل الخارجي، ننشئ \textbf{نسخة جديدة} (مثلاً عبر عملية بنائية ناسخة) في العملية البنائية وفي عمليتي \textenglish{Get} و \textenglish{Set}.

\begin{boxCode}
\begin{english}
\begin{minted}{csharp}
public class Hero
{
    private Point position;

    public Hero(Point start)
    {
        this.position = new Point(start.GetX(), start.GetY()); // Defensive copy
    }

    public Point GetPosition()
    {
        return new Point(this.position.GetX(), this.position.GetY()); // Return copy
    }
}
\end{minted}
\end{english}
\end{boxCode}

\textbf{تنبيه:} المقارنة \textenglish{a == b} بين كائنين تقارن \textbf{العناوين (المؤشرات)} افتراضياً وليس القيم الداخلية. لمقارنة القيم نستخدم عملية مقارنة مثل \textenglish{Equals}.

\filbreak
\subsection*{6. دافع الوراثة: منع تكرار الكود}
\textbf{الفكرة:} إذا اشتركت فئتان أو أكثر في نفس الصفات والعمليات، فإن أي تعديل مستقبلي سيتطلب تعديل عدة أماكن. الحل الهندسي هو استخلاص الجزء \textbf{المتطابق المشترك فقط} في فئة أساسية واحدة، وإبقاء الصفات والعمليات الفريدة في كل فئة فرعية.

\begin{boxCode}
\begin{english}
\begin{minted}{csharp}
// Shared members (name, energy, Eat) belong in base class Animal
// Unique methods (Meow in Cat, Bark in Dog) remain in subclasses
\end{minted}
\end{english}
\end{boxCode}

\filbreak
% ==============================================================
\section*{الجزء الثاني: الوراثة بين الفئات}
% ==============================================================

\subsection*{7. الوراثة واستدعاء العملية البنائية للأب (\textenglish{base})}
\textbf{الفكرة:} كتابة \textenglish{class Dog : Animal} تعني أن \textenglish{Dog} يرث من \textenglish{Animal} جميع صفاته وعملياته المعرّفة بإذن وصول \textenglish{public} أو \textenglish{protected}. \textbf{العمليات البنائية لا تُورَّث}، لذلك تستدعي العملية البنائية للابن العملية البنائية للأب باستخدام \textenglish{: base(...)}.

\begin{boxCode}
\begin{english}
\begin{minted}{csharp}
public class Animal
{
    protected string name;
    protected int energy;

    public Animal(string name, int energy)
    {
        this.name = name;
        this.energy = energy;
    }

    public void Eat(int x) { if (x > 0) { this.energy += x; } }
}

public class Dog : Animal
{
    public Dog(string name) : base(name, 10) { } // Calls Animal constructor
    public void Bark() { Console.WriteLine("Woof"); }
}

// Usage in Main:
Dog d = new Dog("Rex");
d.Eat(5);  // Inherited from Animal
d.Bark();  // Specific to Dog
\end{minted}
\end{english}
\end{boxCode}

\textbf{تنبيه:} لا تعيدوا تعريف الصفات الموروثة (\textenglish{name}, \textenglish{energy}) داخل الفئة الفرعية. إذن الوصول \textenglish{private} في الأب لا يراه الابن؛ نستخدم \textenglish{protected} ليتمكن الأبناء من الوصول إلى الصفات.

\filbreak
\subsection*{8. تمرير كائن كمُعطى لعملية}
\textbf{الفكرة:} عند تمرير كائن إلى عملية، يُمرَّر مؤشره، فأي تعديل عليه داخل العملية يؤثر على الكائن الأصلي. يجب دائماً فحص \textenglish{target != null} والتأكد من تحقق شروط الطرف الأول قبل تعديل الطرف الثاني.

\begin{boxCode}
\begin{english}
\begin{minted}{csharp}
public bool Give(Wallet target, int amount)
{
    if (target == null || amount <= 0 || amount > this.coins)
    {
        return false;
    }
    this.coins -= amount;
    target.Add(amount);
    return true;
}
\end{minted}
\end{english}
\end{boxCode}

\filbreak
% ==============================================================
\section*{الجزء الثالث: تعدد الأشكال}
% ==============================================================

\subsection*{9. مصفوفة الكائنات والتحويل لأعلى}
\textbf{الفكرة:} متغير من نمط الفئة الأساسية (\textenglish{Animal}) يستطيع الإشارة إلى كائن من أي فئة فرعية ترث منها (\textenglish{Dog} أو \textenglish{Cat}). لذلك يمكن لمصفوفة واحدة \textenglish{Animal[]} أن تجمع كائنات من جميع الفئات الفرعية.

\begin{boxCode}
\begin{english}
\begin{minted}{csharp}
public class Zoo
{
    private Animal[] animals;
    public Zoo(int size) { this.animals = new Animal[size]; }

    public bool Add(Animal a)
    {
        if (a == null) { return false; }
        for (int i = 0; i < this.animals.Length; i++)
        {
            if (this.animals[i] == null)
            {
                this.animals[i] = a;
                return true;
            }
        }
        return false;
    }
}

// Usage in Main:
Zoo z = new Zoo(3);
z.Add(new Dog("Rex")); // Upcasting: Dog -> Animal
z.Add(new Cat("Tom")); // Upcasting: Cat -> Animal
\end{minted}
\end{english}
\end{boxCode}

\textbf{تنبيه:} الحدود الفارغة في المصفوفة قيمتها \textenglish{null}؛ لذا افحصوا \textenglish{animals[i] != null} دائماً قبل استدعاء أي عملية عليها لتجنب \textenglish{NullReferenceException}.

\filbreak
\subsection*{10. إعادة تعريف العمليات (\textenglish{virtual} و \textenglish{override})}
\textbf{الفكرة:} العملية التي نريد أن يخصصها كل ابن بطريقته نعرّفها في الأب بالكلمة \textenglish{virtual}، ويعيد الابن تعريفها بالكلمة \textenglish{override} \textbf{بنفس التعريف تماماً}. عند استدعائها عبر مؤشر من نمط الأب، تُنفَّذ نسخة \textbf{الكائن الفعلي} وقت التشغيل.

\begin{boxCode}
\begin{english}
\begin{minted}{csharp}
public class Animal
{
    public virtual string Speak() { return "..."; }
}
public class Dog : Animal
{
    public override string Speak() { return "Woof"; }
}
public class Cat : Animal
{
    public override string Speak() { return "Meow"; }
}

// Usage in Main:
Animal a1 = new Dog("Rex");
Animal a2 = new Cat("Tom");
Console.WriteLine(a1.Speak() + ", " + a2.Speak()); // Prints: Woof, Meow
\end{minted}
\end{english}
\end{boxCode}

\filbreak
\subsection*{11. إرجاع مصفوفة مُصفّاة (مرشّحة) حسب شرط}
\textbf{الفكرة:} لإرجاع مصفوفة جديدة بحجم مطابق تماماً لعدد الكائنات التي تحقق شرطاً معيناً، نمسح المصفوفة مرتين: المرة الأولى لعدّ الحدود المطابقة وإنشاء المصفوفة الجديدة، والمرة الثانية لنسخ المؤشرات إليها.

\begin{boxCode}
\begin{english}
\begin{minted}{csharp}
public Animal[] GetHungry()
{
    int count = 0;
    for (int i = 0; i < this.animals.Length; i++)
    {
        if (this.animals[i] != null && this.animals[i].GetEnergy() < 10) { count++; }
    }
    Animal[] result = new Animal[count];
    int k = 0;
    for (int i = 0; i < this.animals.Length; i++)
    {
        if (this.animals[i] != null && this.animals[i].GetEnergy() < 10)
        {
            result[k] = this.animals[i];
            k++;
        }
    }
    return result;
}
\end{minted}
\end{english}
\end{boxCode}

\filbreak
\subsection*{12. أذونات الوصول \textenglish{protected} والاستدعاء \textenglish{base.Method()}}
\textbf{الفكرة:} داخل عملية مُعاد تعريفها (\textenglish{override}) في الفئة الفرعية، نستدعي نسخة الفئة الأساسية عبر \textenglish{base.Method()} للاستفادة من نتيجتها وإضافة الجزء الخاص بالابن عليها (مثل \textenglish{base.ToString()} أو \textenglish{base.GetSalary()}).

\begin{boxCode}
\begin{english}
\begin{minted}{csharp}
public class Animal
{
    protected string kind;
    protected string name;
    public override string ToString() { return $"Kind:{this.kind}, Name:{this.name}"; }
}

public class Dog : Animal
{
    private string breed;
    public override string ToString() { return base.ToString() + $", Breed:{this.breed}"; }
}
\end{minted}
\end{english}
\end{boxCode}

\filbreak
\subsection*{13. النمط الساكن للمتغير مقابل النمط الفعلي للكائن}
\textbf{الفكرة:} المترجم يفحص قانونية استدعاء العملية حسب \textbf{نمط المتغير (الساكن)}، بينما وقت التشغيل ينفّذ العملية حسب \textbf{نمط الكائن (الفعلي)}. لذلك لا يمكن عبر متغير من نمط \textenglish{Animal} استدعاء عملية موجودة في \textenglish{Dog} فقط إلا بعد إجراء تحويل لأسفل.

\begin{boxCode}
\begin{english}
\begin{minted}{csharp}
Animal a = new Dog("Rex");
a.Speak();        // Valid: Speak() is declared in Animal
// a.Bark();      // Compile error: Animal does not have Bark()
((Dog)a).Bark();  // Valid after downcasting to Dog
\end{minted}
\end{english}
\end{boxCode}

\filbreak
\subsection*{14. العامل \textenglish{is} والتحويل لأسفل}
\textbf{الفكرة:} التعبير البولياني \textenglish{x is Dog} يفحص وقت التشغيل إن كان الكائن الفعلي من نمط \textenglish{Dog} (ويعيد \textenglish{false} بأمان إذا كان \textenglish{null}). بعد التأكد من نمطه، نحوله بالصيغة \textenglish{(Dog)x} لاستدعاء عملياته الخاصة.

\begin{boxCode}
\begin{english}
\begin{minted}{csharp}
for (int i = 0; i < this.animals.Length; i++)
{
    if (this.animals[i] is Dog)
    {
        Dog d = (Dog)this.animals[i];
        d.Bark();
    }
}
\end{minted}
\end{english}
\end{boxCode}

\textbf{تنبيه:} إجراء تحويل مباشر \textenglish{(Dog)a} دون فحص مسبق بـ \textenglish{is} يسبب انهيار البرنامج (\textenglish{InvalidCastException}) إذا كان الكائن الفعلي من نمط آخر.

\filbreak
\subsection*{15. العامل \textenglish{as} ومقارنة أساليب التحويل}
\textbf{الفكرة:} التعبير \textenglish{x as Dog} يحاول تحويل الكائن إلى \textenglish{Dog}؛ فإذا نجح يعيد مؤشر الكائن، وإذا فشل يعيد \textenglish{null} دون رمي استثناء. لذلك نفحص بعده \textenglish{d != null}.

\begin{boxCode}
\begin{english}
\begin{minted}{csharp}
Dog d = this.animals[i] as Dog;
if (d != null)
{
    d.Bark();
}
\end{minted}
\end{english}
\end{boxCode}

\begin{center}
\renewcommand{\arraystretch}{1.2}
{\small
\begin{tabular}{|c|c|c|c|}
\hline
\textbf{الأسلوب} & \textbf{عند التوافق} & \textbf{عند عدم التوافق} & \textbf{مع \textenglish{null}} \\
\hline
\textenglish{is} ثم \textenglish{(Dog)x} & ينجح & يمنع التحويل بأمان & \textenglish{false} \\
\hline
\textenglish{x as Dog} & ينجح & يعيد \textenglish{null} & يعيد \textenglish{null} \\
\hline
\textenglish{(Dog)x} مباشرة & ينجح & \textbf{انهيار (\textenglish{InvalidCastException})} & \textenglish{null} \\
\hline
\end{tabular}
}
\end{center}

\filbreak
\subsection*{16. قوة تعدد الأشكال وقابلية التوسع}
\textbf{الفكرة:} عندما تتعامل العمليات في فئة الإدارة مع الواجهة المشتركة للفئة الأساسية (مثل \textenglish{Speak()} أو \textenglish{Eat()})، فإن إضافة فئة فرعية جديدة مستقبلاً (مثل \textenglish{Bird}) تعمل تلقائياً دون تعديل كود فئة الإدارة.

\begin{boxCode}
\begin{english}
\begin{minted}{csharp}
public void MakeAllSpeak()
{
    for (int i = 0; i < this.animals.Length; i++)
    {
        if (this.animals[i] != null) { Console.WriteLine(this.animals[i].Speak()); }
    }
}
\end{minted}
\end{english}
\end{boxCode}

\filbreak
\subsection*{17. الوراثة من الفئة \textenglish{Object} — العملية \textenglish{Equals}}
\textbf{الفكرة:} جميع الفئات في \textenglish{C\#} ترث ضمنياً من الفئة العليا \textenglish{object}. لمقارنة كائنين حسب قيم صفاتهما، نعيد تعريف العملية \textenglish{Equals(object obj)} بالترتيب: فحص \textenglish{null} ونمط الكائن عبر \textenglish{is}، ثم التحويل، ثم مقارنة الصفات.

\begin{boxCode}
\begin{english}
\begin{minted}{csharp}
public override bool Equals(object obj)
{
    if (obj == null || !(obj is Animal)) { return false; }
    Animal other = (Animal)obj;
    return this.id == other.id;
}
\end{minted}
\end{english}
\end{boxCode}

\filbreak
% ==============================================================
\section*{الجزء الرابع: الفئات المجردة والواجهات}
% ==============================================================

\subsection*{18. الفئة المجردة (\textenglish{abstract}) والعمليات المجردة}
\textbf{الفكرة:} الفئة المجردة هي فئة أساسية غير مكتملة لا يمكن إنشاء كائنات منها مباشرة (\textenglish{new Animal()} ممنوع)، وتُستخدم كفئة أب فقط. العملية المجردة (\textenglish{abstract}) تُكتب بتعريفها فقط دون جسم (تنتهي بـ \textenglish{;})، ويلتزم كل ابن غير مجرد بإعادة تعريفها (\textenglish{override}).

\begin{boxCode}
\begin{english}
\begin{minted}{csharp}
public abstract class Animal
{
    protected string name;
    public abstract string Speak(); // No body - must be overridden
}

public class Dog : Animal
{
    public override string Speak() { return "Woof"; }
}
\end{minted}
\end{english}
\end{boxCode}

\filbreak
\subsection*{19. المُعطى العام من نمط \textenglish{object} وفحص نمطه}
\textbf{الفكرة:} عندما تشترك الفئات الوارثة في عملية مجردة واحدة لكن تختلف أنماط المعطيات التي تحتاجها كل فئة، نجعل المُعطى من نمط \textenglish{object}، وداخل كل فئة فرعية نفحص نمطه بـ \textenglish{is} ونحوله قبل الاستخدام.

\begin{boxCode}
\begin{english}
\begin{minted}{csharp}
public class Dog : Animal
{
    public override bool Feed(object o)
    {
        if (o is int)
        {
            int amount = (int)o;
            if (amount > 0) { this.energy += amount; return true; }
        }
        return false;
    }
}
\end{minted}
\end{english}
\end{boxCode}

\filbreak
\subsection*{20. الواجهة \textenglish{IComparable} والعملية \textenglish{CompareTo}}
\textbf{الفكرة:} الواجهة هي عقد يحدد عمليات يجب على الفئة تطبيقها. الواجهة \textenglish{IComparable} تلزم الفئة بتطبيق العملية \textenglish{int CompareTo(object obj)} التي تقارن الكائن الحالي (\textenglish{this}) بالكائن المُتلقّى (\textenglish{obj}) وتعيد: \textbf{عدداً سالباً} إذا كان الكائن الحالي أصغر، \textbf{صفراً} إذا تساويا، و\textbf{عدداً موجباً} إذا كان الكائن الحالي أكبر.

\begin{boxCode}
\begin{english}
\begin{minted}{csharp}
public class Player : IComparable
{
    private int score;
    public int CompareTo(object obj)
    {
        if (obj == null || !(obj is Player)) { return 1; }
        Player other = (Player)obj;
        return this.score - other.score;
    }
}
\end{minted}
\end{english}
\end{boxCode}

\filbreak
\subsection*{21. خدمات الفرز العامة عبر الواجهة \textenglish{IComparable}}
\textbf{الفكرة:} نكتب عملية فرز واحدة تتلقى مصفوفة \textenglish{IComparable[]} وتعتمد على \textenglish{CompareTo}، فتصلح لفرز أي مصفوفة كائنات تطبّق الواجهة \textenglish{IComparable}.

\begin{boxCode}
\begin{english}
\begin{minted}{csharp}
public class SortServices
{
    public static void SortAscending(IComparable[] arr)
    {
        for (int i = 0; i < arr.Length - 1; i++)
        {
            for (int j = 0; j < arr.Length - 1 - i; j++)
            {
                if (arr[j].CompareTo(arr[j + 1]) > 0)
                {
                    IComparable temp = arr[j];
                    arr[j] = arr[j + 1];
                    arr[j + 1] = temp;
                }
            }
        }
    }
}
\end{minted}
\end{english}
\end{boxCode}

\filbreak
\subsection*{22. الدمج الشامل: تصفية حسب النمط ثم فرز وإرجاع مصفوفة}
\textbf{الفكرة:} لاستخراج مصفوفة مرتّبة لكائنات من فئة فرعية معينة: (1) نمسح المصفوفة ونعدّ الكائنات من النمط المطلوب باستخدام \textenglish{is}، (2) ننسخ مؤشراتها إلى مصفوفة جديدة بحجمها الدقيق، (3) نفرز المصفوفة الجديدة عبر \textenglish{SortServices} ونعيدها دون تغيير ترتيب المصفوفة الأصلية.

\begin{boxCode}
\begin{english}
\begin{minted}{csharp}
public Dog[] GetDogsSorted()
{
    int count = 0;
    for (int i = 0; i < this.animals.Length; i++)
    {
        if (this.animals[i] is Dog) { count++; }
    }
    Dog[] result = new Dog[count];
    int k = 0;
    for (int i = 0; i < this.animals.Length; i++)
    {
        if (this.animals[i] is Dog) { result[k] = (Dog)this.animals[i]; k++; }
    }
    SortServices.SortAscending(result);
    return result;
}
\end{minted}
\end{english}
\end{boxCode}

\filbreak
\begin{boxAttention}[- قائمة مراجعة سريعة قبل تسليم أي تمرين]
\begin{enumerate}[itemsep=0.2em]
    \item هل كتبتُ في \textenglish{Main} كوداً يختبر كل عملية جديدة بحالات ناجحة وأخرى فاشلة؟
    \item هل فحصتُ \textenglish{!= null} قبل استخدام أي حد في مصفوفة أو مُعطى من نمط كائن؟
    \item هل كل عملية تحتمل الفشل تفحص جميع الشروط \textbf{قبل} تعديل صفات الكائن؟
    \item هل وضعتُ الكلمات \textenglish{protected}، \textenglish{base}، \textenglish{virtual}، \textenglish{override}، و \textenglish{abstract} في أماكنها الصحيحة؟
    \item هل استعنتُ بـ \textenglish{base.ToString()} في الفئات الفرعية لمنع تكرار الكود؟
\end{enumerate}
\end{boxAttention}

\begin{flushleft}
أرجو لكم وقتًا ممتعًا.

الأستاذ محمود اغبارية.
\end{flushleft}

\end{document}
